%!TEX program = xelatex
%!TEX root = geometria_analitica.tex
%%Usar makeindex -s indexstyle.ist arquivo.idx no terminal para gerar o índice remissivo agrupado por inicial
%%Após executar pdflatex arquivo

\chapter{Circunferências e Cônicas} % (fold)
\label{cha:circunferencias_e_conicas}

\section{Equações da Circunferência} % (fold)
\label{sec:equacoes_da_circunferencias}

\begin{definicao}
	A circunferência é o conjunto dos pontos do plano que estão a uma mesma distância (denominada \textbf{raio}) de um dado ponto do plano (chamado \textbf{centro}).\index{Circunferência}
\end{definicao}

Considere uma circunferência de centro $C(x_0,y_0)$ e raio $r$.  Considere também os pontos $A(x_0 + r, y_0)$ e $B(x_0, y_0 + r)$. Seja $P(x,y)$ um ponto qualquer desta circunferência e seja $\beta$ o ângulo formado pelos vetores $\vec{CA}$ e $\vec{CP}$, onde $A(x_0 + r, y_0)$. Assim
\[
	\cos\beta = \dfrac{\inner{CA}{CP}}{\norma{\vec{CA}}\norm{\vec{CP}}}.
\]

Mas $\vec{CP} = (x - x_0, y - y_0)$, $\vec{CA} = (r, 0)$ e $\norm{\vec{CA}} = \norm{\vec{CP}} = r$, então
\[
	\cos\beta = \dfrac{(x - x_0)r}{r^2} = \dfrac{x - x_0}{r},
\]
isto é,
\begin{equation}\label{equacaoparametricaX}
	x = x_0 + r\cos\beta
\end{equation}

Agora, para os vetores $\vec{CP}$ e $\vec{CB}$, temos $\vec{CP} = (x - x_0, y - y_0)$, $\vec{CB} = (0, r)$ e $\norm{\vec{CA}} = \norm{\vec{CB}} = r$, então
\begin{equation}\label{equacaoinicialY}
	\cos\alpha = \dfrac{y - y_0}{r}.
\end{equation}

\begin{figure}[h]
	\centering
	\caption{Circunferência}\label{circunferencia}
	\includegraphics{circunferencia.pdf}
\end{figure}

Na Figura \ref{circunferencia}, considere o pé da perpendicular ao segmento $CA$ passando por $P$ e denote esse ponto por $E$. Agora, considere o pé da perpendicular ao segmento $CB$ passando por $P$ e denote esse ponto por $D$. Então os segmentos $CD$ e $EP$ têm o mesmo comprimento $k$. Daí
\begin{equation}\label{equacaosenocosseno}
	\cos\alpha = \dfrac{k}{r} = \sin\beta.
\end{equation}

Substituindo \eqref{equacaosenocosseno} em \eqref{equacaoinicialY} obtemos:
\[
	\sin\beta = \dfrac{y - y_0}{r},
\]
isto é,
\begin{equation}\label{equacaoparametricaY}
	y = y_0 + r\sin\beta.
\end{equation}

Por outro lado, se um ponto $P(x, y)$ satisfaz as equações \eqref{equacaoparametricaX} e \eqref{equacaoparametricaY}, então
\[
	d(P,C) = \sqrt{(x - x_0)^2 + (y - y_0)^2} = \sqrt{r^2\cos^2\beta + r^2\sin^2\beta} = r
\]
e então o ponto $P$ pertence á circunferência de centro $C(x_0,y_0)$ e raio $r$.

\begin{definicao}
	As equações
	\[
		\begin{cases}
			x = x_0 + r\cos\beta \\
			y = y_0 + r\sin\beta,
		\end{cases}
		\beta \in [0,2\pi]
	\]
	são chamadas \textbf{equações paramétricas} da circunferência de centro $C(x_0, y_0)$ e raio $r$.\index{Circunferência!Equações Paramétricas}
\end{definicao}

Seja $C$ uma circunferência com equações paramétricas
\[
	\begin{cases}
		x = x_0 + r\cos\beta \\
		y = y_0 + r\sin\beta,
	\end{cases}
	\beta \in [0,2\pi].
\]

Podemos escrever
\begin{align*}
	x - x_0 = r\cos\beta \\
	y - y_0 = r\sin\beta.
\end{align*}
Elevando ao quadrado ambos os lados
\begin{align*}
	(x - x_0)^2 = r^2\cos^2\beta \\
	(y - y_0)^2 = r^2\sin^2\beta
\end{align*}
e somando obtemos
\begin{equation}
	(x - x_0)^2 + (y - y_0)^2 = r^2
\end{equation}
que é a \textbf{equação cartesiana} da circunferência de centro $(x_0,y_0)$ e raio $r$.\index{Circunferência!Equação Cartesiana}
Podemos reescrever essa equação como
\[
	x^2 + y^2 - 2x_0x - 2y_0y + x_0^2 + y_0^2 - r^2 = 0.
\]
Fazendo
\[
	-2x_0 = a,\quad -2y_0 = b,\quad x_0^2 + y_0^2 - r^2 = c
\]
obtemos
\begin{equation}\label{equacaoalternativacircunferencia}
	x^2 + y^2 + ax + by + c = 0.
\end{equation}

Sejam $P_1(x_1,y_1)$, $P_2(x_2,y_2)$ e $P_3(x_3,y_3)$ três pontos não colineares, isto é, os vetores $\vec{P_1P_2}$ e $\vec{P_1P_3}$ não são paralelos. Nessa situação é possível obter uma circunferência contendo estes três pontos. Para isso considere a equação \eqref{equacaoalternativacircunferencia}. Para que $P_1$, $P_2$ e $P_3$ pertençam a uma circunferência $C$ devemos ter:
\begin{align*}
	x_1^2 + y_1^2 + ax_1 + by_1 + c = 0 \\
	x_2^2 + y_2^2 + ax_2 + by_2 + c = 0 \\
	x_3^2 + y_3^2 + ax_3 + by_3 + c = 0.
\end{align*}
Daí para encontrar a circunferência desejada, basta determinar os valores de $a$, $b$ e $c$. Ou seja, precisamos resolver o sistema
\[
	\begin{cases}
		ax_1 + by_1 + c = -(x_1^2 + y_1^2) \\
		ax_2 + by_2 + c = -(x_2^2 + y_2^2) \\
		ax_3 + by_3 + c = -(x_3^2 + y_3^2).
	\end{cases}
\]

Como os ponto $P_1$, $P_2$ e $P_3$ não são colineares, este sistema sempre possui solução única.

\begin{exemplos}
	\begin{enumerate}
		\item Escreva a equação cartesiana da circunferência passando pelos pontos $A(1,1)$, $B(1,-2)$ e $C(2,3)$.
		      \begin{solucao}
			      Temos
			      \begin{align*}
				      \vec{AB} = (0,-3) \\
				      \vec{AC} = (1,2)
			      \end{align*}
			      assim os pontos não são colineares e portanto definem uma circunferência. Para encontrá-la, devemos resolver o sistema
			      \[
				      \begin{cases}
					      a + b + c = -2  \\
					      a - 2b + c = -5 \\
					      2a + 3b + c = -13
				      \end{cases}
			      \]
			      cuja solução é
			      \[
				      a = -13,\quad b = 1 \quad\mbox{e}\quad c = 10.
			      \]
			      Logo
			      \[
				      x^2 + y^2 - 13x + y + 10 = 0
			      \]
			      é a equação procurada.

			      Agora, como
			      \[
				      -2x_0 = a,\quad -2y_0 = b,\quad x_0^2 + y_0^2 - r^2 = c
			      \]
			      segue que
			      \[
				      x_0 = 13/2,\quad y_0 = -1/2 \quad\mbox{e}\quad r^2 = 65/2.
			      \]
			      Assim a equação cartesiana da circunferência será
			      \[
				      \left(x - \dfrac{13}{2}\right)^2 + \left(y + \dfrac{1}{2}\right)^2 = \dfrac{65}{2}.
			      \]
			      As equações paramétricas são
			      \[
				      \begin{cases}
					      x = \dfrac{13}{2} + \dfrac{\sqrt{65}}{\sqrt{2}}\cos\beta \\
					      y = -\dfrac{1}{2} + \dfrac{\sqrt{65}}{\sqrt{2}}\sin\beta \\
				      \end{cases}, \beta \in [0,2\pi].
			      \]
		      \end{solucao}
		\item Encontre a equação da circunferência que passa pelos pontos $P_1(1,-1)$, $P_2(0,1)$ e $P_3(1,0)$.
		      \begin{solucao}
			      Primeiro note que
			      \begin{align*}
				      \vec{P_1P_2} = (-1,2) \\
				      \vec{P_1P_3} = (0,1).
			      \end{align*}
			      Assim os pontos não são colineares, portanto tal circunferência existe. Para encontrá-la, basta resolver o sistema
			      \[
				      \begin{cases}
					      a - b + c = -2 \\
					      b + c = -1     \\
					      a + c = -2.
				      \end{cases}
			      \]
			      A solução é $a = b = 1$ e $c = -2$. Logo a circunferência tem equação
			      \[
				      x^2 + y^2 + x + y - 2 = 0.
			      \]
		      \end{solucao}
		\item Encontre a equação da circunferência que passa pelos pontos $P_1(1,3)$, $P_2(3,-5)$ e $P_3(2,-1)$.
		      \begin{solucao}
			      Note que
			      \begin{align*}
				      \vec{P_1P_2} = (2, -8) \\
				      \vec{P_1P_3} = (1, -4).
			      \end{align*}
			      Assim $P_1$, $P_2$ e $P_3$ são colineares e portanto os três pontos não definem uma circunferência.
		      \end{solucao}
	\end{enumerate}
\end{exemplos}

% section equacoes_da_circunferencias (end)

\section{Cônicas} % (fold)
\label{sec:conicas}

\begin{definicao}
	Chama-se \textbf{cônica} o lugar geométrico dos pontos $P(x,y)$ que satisfazem uma equação do segundo grau $g(x,y) = 0$, onde\index{Cônicas}
	\begin{equation}\label{equacaoconica}
		g(x,y) = ax^2 + bxy + cy^2 + dx + ey + f.
	\end{equation}
\end{definicao}

Na equação \eqref{equacaoconica} pelo menos um dos números $a$, $b$ $c$ é diferente de zero. Os termos $ax^2$ e $cy^2$ são os \textbf{termos quadráticos}, o termo $bxy$ é o \textbf{termo quadrático misto}, $dx$ e $ey$ são os \textbf{termos lineares} e $f$ é o \textbf{termo independente}. \index{Cônica!Termos Quadráticos} \index{Cônica!Termo Quadrático Misto} \index{Cônica!Termos Lineares} \index{Cônica!Termos Independete}

A equação \eqref{equacaoconica} pode representar:
\begin{enumerate}[label=({\arabic*})]
	\item um conjunto vazio: $x^2 + y^2 = -1$;
	\item um conjunto formado por um ponto: $x^2 + y^2 = 0$;
	      %\item uma reta;
	\item a reunião de duas retas (paralelas ou concorrentes): $x^2 - y^2 - 2x - 6y -8 =0$;
	\item uma circunferência: $x^2 + y^2 - 13x + y + 10 = 0$;
	\item uma elipse: $4x^2 + 9y^2 - 8x - 36y + 4 = 0$;
	\item uma hipérbole: $3x^2 + 3y^2 - 10xy + 12\sqrt{2}x - 4\sqrt{2}y + 32 = 0$;
	\item uma parábola: $x^2 - 2xy + y^2 -2x - 2y + 1 = 0$.
\end{enumerate}

Nos casos em que a equação \eqref{equacaoconica} representar um conjunto vazio, um ponto, retas ou circunferências diremos que a cônica definida é uma \textbf{cônica degenerada}. Aqui vamos estudar as cônicas chamadas de elipse, hipérbole e parábola.

\subsection{Elipse} % (fold)
\label{sub:elipse}

\begin{definicao}
	Sejam $F_1$ e $F_2$ pontos distintos do plano, $2c$ a distância entre $F_1$ e $F_2$ e $a$ um número real tal que $a > c$. O lugar geométrico $\mathcal{E}$ dos pontos $P$ tais que
	\begin{equation}\label{definicaoelipse}
		d(P,F_1) + d(P,F_2) = 2a
	\end{equation}
	chama-se \textbf{elipse}. Cada um dos pontos $F_1$ e $F_2$ é chamado \textbf{foco} da elipse. O segmento $F_1F_2$ é chamado \textbf{segmento focal}, seu ponto médio é o \textbf{centro} da elipse e $2c$ é chamado \textbf{distância focal}. A reta passando por $F_1$ e $F_2$ chama-se \textbf{reta focal}.\index{Elipse} \index{Elipse!Focos} \index{Elipse!Centro} \index{Elipse!Reta Focal}
\end{definicao}

Para obter a equação que descreve uma elipse vamos fixar os focos $F_1$ e $F_2$ como sendo $F_1(-c,0)$ e $F_2(c,0)$. Com isso o centro dessa elipse estará na
origem do plano cartesiano. Assim, seja $P(x,y)$ um ponto pertencente à elipse de focos $F_1$ e $F_2$. Então
\begin{equation}
	d(P,F_1) + d(P,F_2) = 2a,
\end{equation}
isto é,
\begin{equation}
	d(P,F_1) = 2a -  d(P,F_2).
\end{equation}
Assim,
\[
	\sqrt{(x + c)^2 + y^2} = 2a - \sqrt{(x - c)^2 + y^2}
\]
e elevando ao quadrado ambos os membros dessa igualdade obtemos
\[
	a\sqrt{(x - c)^2 + y^2} = a^2 - cx.
\]
Elevando novamente ao quadrado e simplificando
\[
	(a^2 - c^2)x^2 + a^2y^2 = a^2(a^2 - c^2).
\]
Fazendo $b^2 = a^2 - c^2$, podemos escrever
\[
	b^2x^2 + a^2y^2 = a^2b^2.
\]
Finalmente, dividindo por $a^2b^2$ obtemos
\begin{equation}\label{equacaoelipse}
	\elipse{a^2}{b^2}.
\end{equation}

Observe que
\begin{align}
	a^2 = b^2 + c^2 \\
	a > b > 0       \\
	a > c > 0.
\end{align}

Assim todo ponto pertencente a elipse de focos $F_1(-c,0)$ e $F_2(c,0)$, deve satisfazer a equação
\[
	\elipse{a^2}{b^2}.
\]

Por outro lado, seja $Q(\alpha,\beta)$ um ponto tal que
\[
	\dfrac{\alpha^2}{a^2} + \dfrac{\beta^2}{b^2} = 1.
\]
Então $\beta^2 = b^2 - (b^2/a^2)\alpha^2$ e daí
\begin{align*}
	d^2(Q,F_1) & = (\alpha + c)^2 + \beta^2                                   \\
	           & = \alpha^2 + 2c\alpha + c^2 + b^2 - \dfrac{b^2}{a^2}\alpha^2 \\
	           & = \left(\dfrac{c}{a}\alpha + a\right)^2.
\end{align*}

Analogamente,
\begin{equation*}
	d^2(Q,F_2) = \left(\dfrac{c}{a}\alpha - a\right)^2.
\end{equation*}
Assim
\begin{align}
	d(Q,F_1) & = \left|\dfrac{c}{a}\alpha + a\right|\label{distanciaP1}  \\
	d(Q,F_2) & = \left|\dfrac{c}{a}\alpha - a\right|\label{distanciaP2}.
\end{align}

Como
\[
	\dfrac{\alpha^2}{a^2} + \dfrac{\beta^2}{b^2} = 1.
\]
então $\alpha^2/a^2 \le 1$, isto é, $\mid\alpha\mid\le a$. Mas $c/a > 0$, daí
\[
	\dfrac{c}{a}\mid\alpha\mid \le c < a
\]
logo
\[
	-a < \dfrac{c\alpha}{a} < a.
\]
Portanto,
\begin{align*}
	\dfrac{c}{a}\alpha + a > 0 \\
	\dfrac{c}{a}\alpha - a < 0.
\end{align*}
Desse modo podemos escrever \eqref{distanciaP1} e \eqref{distanciaP2} como
\begin{align*}
	d^2(Q,F_1) & = \dfrac{c}{a}\alpha + a  \\
	d^2(Q,F_2) & = a - \dfrac{c}{a}\alpha.
\end{align*}

Assim
\[
	d(Q,F_1) + d(Q,F_2) = 2a
\]
e então o ponto $Q(\alpha,\beta)$ pertence á elipse de focos $F_1(-c,0)$ e $F_2(c,0)$.

Com isso provamos a seguinte proposição:
\begin{proposicao}
	Um ponto $P(x,y)$ é um ponto da elipse de equação reduzida
	\[
		\elipse{a^2}{b^2}
	\]
	se, e somente se, as distâncias de $P$ aos focos $F_1$ e $F_2$ são
	\begin{align*}
		d^2(P,F_1) & = \dfrac{c}{a}\alpha + a  \\
		d^2(P,F_2) & = a - \dfrac{c}{a}\alpha.
	\end{align*}
\end{proposicao}


Os números $a$, $b$ e $c$ são chamados \textbf{parâmetros geométricos} da elipse. A equação \eqref{equacaoelipse} é chamada de \textbf{equação reduzida da elipse} de centro $O(0,0)$ e focos no eixo $x$. Indicamos tal elipse por\index{Elipse!Parâmetros Geométricos} \index{Elipse!Equação Reduzida}
\[
	\mathcal{E}: \elipse{a^2}{b^2}.
\]

\begin{observacao}
	\begin{enumerate}
		\item A equação \eqref{equacaoelipse} só é válida para elipses com focos no eixo $x$ e centro na origem. Mudanças nos focos ou no centro, resultarão em uma equação diferente para a elipse.
		\item Se $(x,y)$ satisfaz a equação \eqref{equacaoelipse}, então $(-x,-y)$, $(-x,y)$ e $(x,-y)$ também satisfazem a mesma equação. Assim a elipse é simétrica em relação ao eixo $x$, ao eixo $y$ e ao centro do plano cartesiano.
		\item O ponto $(x,0)$ pertence à elipse de equação \eqref{equacaoelipse} se, e somente se, $x^2/a^2 = 1$, isto é, $x = \pm a$. Além disso, o ponto
		      $(0,y)$ pertence à elipse de equação \eqref{equacaoelipse} se, e só se, $y = \pm b$. Portanto, a interseção de $\mathcal{E}$ com o eixo $x$ são os pontos $A_1(-a,0)$ e $A_2(a,0)$ e a interseção com o eixo $y$ são os pontos $B_1(0,-b)$ e $B_2(0,b)$. Note que $d(B_j,F_i) = \sqrt{c^2 + b^2} = a$, para $i$, $j = 1$, 2.
		\item Não existe circunferência que contenha os pontos $A_1$, $A_2$, $B_1$ e $B_2$. Caso existisse, tal circunferência teria dois diâmetros diferentes,
		      cujos comprimento são dados pelos comprimentos dos segmentos $A_1A_2$ e $B_1B_2$, o que é impossível. Portanto, como esses pontos pertencem a
		      $\mathcal{E} $, podemos afirmar que a elipse não é uma circunferência e nem o conjunto vazio.
	\end{enumerate}
\end{observacao}

Na equação \eqref{equacaoelipse}, fazendo $y \ge 0$ podemos escrever
\[
	y = \dfrac{b}{a}\sqrt{a^2 - x^2},\quad x \in [0,a].
\]
Tal equação é uma função em $x$, contínua, decrescente e com concavidade para baixo em todo seu domínio. Assim seu gráfico é dado pela Figura \ref{GraficoprimeiroquadranteElipse}.
\begin{figure}[h!]
	\centering
	\caption{Gráfico da função $y = \dfrac{b}{a}\sqrt{a^2 - x^2}$}
	\label{GraficoprimeiroquadranteElipse}
	\includegraphics{grafico-elipse-primeiro-quadrante.pdf}
\end{figure}

E como a elipse é simétrica, sua forma é dada pela Figura \ref{FormageralElipse}.
\begin{figure}[h]
	\centering
	\caption{Elipse com focos no eixo $x$ e centro na origem: $\elipse{a^2}{b^2}$}
	\label{FormageralElipse}
	\includegraphics{elipse-focos-eixo-X.pdf}
\end{figure}

\begin{definicao}
	Os pontos $A_1$ e $A_2$ em que a reta focal intercepta a elipse e os pontos $B_1$ e $B_2$ em que a mediatriz do segmento focal intercepta a elipse são chamados de \textbf{vértices}. Os segmentos $A_1A_2$ e $B_1B_2$ são chamados de \textbf{eixo maior} e \textbf{eixo menor} da elipse, respectivamente.\index{Elipse!Vértice} \index{Elipse!Eixo Maior} \index{Elipse!Eixo Menor}
\end{definicao}

\begin{exemplos}
	\begin{enumerate}
		\item A equação
		      \[
			      \elipse{9}{4}
		      \]
		      representa uma elipse onde $a = 3$, $b = 2$. Os focos são $F_1(-c,0)$ e $F_2(c,0)$ onde, $a^2 = b^2 + c^2$. Daí, $c = \sqrt{5}$. Logo $F_1(-\sqrt{5},0)$ e $F_2(\sqrt{5},0)$. Os vértices são $A_1(-3,0)$, $A_2(3,0)$, $B_1(0,-2)$ e $B_2(0,2)$.
		\item Encontre a equação da elipse cujos focos são $F_1(-3,0)$ e $F_2(0,4)$ e cujo eixo maior mede 7.
		      \begin{solucao}
			      Como os focos não estão sobre o eixo $x$, não podemos usar a equação \eqref{equacaoelipse}. Assim vamos utilizar a definição \eqref{definicaoelipse}. Seja $P(x,y)$ um ponto pertencente a tal elipse, então
			      \[
				      d(P,F_1) + d(P,F_2) = 7.
			      \]
			      Assim
			      \begin{align*}
				       & \sqrt{(x + 3)^2 + y^2} + \sqrt{x^2 + (y - 4)^2} = 7                  \\
				       & (x+ 3)^2 + y^2 = 49 - 14\sqrt{x^2 + (y - 4)^2} + x^2 + (y - 4)^2     \\
				       & (3x + 4y - 28)^2 = \left[-7\sqrt{x^2 + (y - 4)^2}\right]^2           \\
				       & 9x^2 + 24xy + 16y^2 - 168x - 224y + 784 = 49x^2 + 49y^2 - 392y + 784
			      \end{align*}
			      Portanto a elipse procurada tem equação
			      \[
				      40x^2 + 33y^2 - 24xy + 168x - 168y = 0.
			      \]
		      \end{solucao}
		\item Encontre a equação da elipse com focos $F_1(0,-2)$ e $F_2(0,2)$ e tal que $a = 3$.
		      \begin{solucao}
			      Neste caso, $P(x,y)$ pertence à elipse se, e somente se,
			      \begin{align*}
				       & d(P,F_1) + d(P,F_2) = 6                                \\
				       & \sqrt{x^2 + (y + 2)^2} + \sqrt{x^2 + (y - 2)^2} = 6    \\
				       & (8y - 36)^2 = \left[-12\sqrt{x^2 + (y - 2)^2}\right]^2 \\
				       & 64y^2 - 576y + 1296 = 144x^2 + 144y^2 - 576y + 576     \\
				       & 144x^2 + 80y^2 = 720.
			      \end{align*}
			      Portanto a elipse procurada tem equação
			      \[
				      \elipse{5}{9}.
			      \]
		      \end{solucao}
	\end{enumerate}
\end{exemplos}

\begin{proposicao}
	Uma equação da forma
	\begin{equation}\label{equacaoelipsequasegeral}
		\elipse{p}{q}
	\end{equation}
	descreve uma elipse se, e somente se, os números reais $p$ e $q$ são distintos e positivos.
\end{proposicao}

\begin{corolario}
	Sejam $O$ a origem do plano cartesiano, $p$ e $q$ números reais distintos e positivos. Seja $\mathcal{E}$ a elipse de equação \eqref{equacaoelipsequasegeral} e parâmetros geométricos $a$ e $b$.
	\begin{enumerate}
		\item Se $p > q$, então $a^2 = p$, $b^2 = q$ e $\mathcal{E}$ tem centro $O$ e focos no eixo $x$.
		\item Se $q > p$, então $a^2 = q$, $b^2 = p$ e $\mathcal{E}$ tem centro $O$ e focos no eixo $y$.
	\end{enumerate}
\end{corolario}

\begin{figure}[h]
	\centering
	\caption{Elipse com focos no eixo $y$ e centro na origem: $\elipse{a^2}{b^2}$}
	\includegraphics{elipse-focos-eixo-Y.pdf}
\end{figure}


\begin{exemplos}
	Para as equações dadas, encontre os focos, os vértices e o comprimento do eixo maior e menor da elipse dada.
	\begin{enumerate}
		\item $16x^2 + y^2 = 1$
		      \begin{solucao}
			      Podemos escrever essa equação como
			      \[
				      \elipse{\dfrac{1}{16}}{1}
			      \]
			      e daí $p = 1/16$ e $q = 1$. Como $p < q$, os focos estão no eixo $y$. Assim $a = 1$ e $b = 1/4$. Como $a^2 = b^2 + c^2$, temos $c = \sqrt{15}/4$.

			      Os focos são:
			      \[
				      F_1\left(0,-\dfrac{\sqrt{15}}{4}\right),\quad F_2\left(0,\dfrac{\sqrt{15}}{4}\right).
			      \]

			      Os vértices são:
			      \begin{align*}
				      A_1(0,-1),\quad A_2(0,1) \\
				      B_1\left(-\dfrac{1}{4},0\right),\quad B_1\left(\dfrac{1}{4},0\right).
			      \end{align*}
			      O eixo maior tem comprimento 2 e o eixo menor tem comprimento 1/2.
		      \end{solucao}
		\item $25x^2 + 169y^2 = 9$
		      \begin{solucao}
			      Podemos escrever essa equação como
			      \[
				      \elipse{\dfrac{9}{25}}{\dfrac{9}{169}}
			      \]
			      e daí $p = 9/25$ e $q = 9/169$. Como $p > q$, os focos estão no eixo $x$. Assim $a = 3/5$ e $b = 3/13$. Como $a^2 = b^2 + c^2$, temos $c = 36/65$.

			      Os focos são:
			      \[
				      F_1\left(-\dfrac{36}{65}\right),\quad F_2\left(\dfrac{36}{65}\right).
			      \]

			      Os vértice são:

			      \begin{align*}
				      A_1\left(-\dfrac{3}{5},0\right),\quad A_2\left(\dfrac{3}{5},0\right) \\
				      B_1\left(0,-\dfrac{3}{13}\right),\quad B_1\left(0,\dfrac{3}{13}\right).
			      \end{align*}
			      O eixo maior tem comprimento 1,2 e o eixo menor tem comprimento 6/13.
		      \end{solucao}
	\end{enumerate}
\end{exemplos}

Agora considere a elipse de equação reduzida
\[
	\mathcal{E}: \elipse{a^2}{b^2},\quad a > b.
\]
Temos
\begin{align*}
	\dfrac{x^2}{a^2} \le 1 \Leftrightarrow -a \le x \le a \\
	\dfrac{y^2}{b^2} \le 1 \Leftrightarrow -b \le x \le b.
\end{align*}

Portanto os pontos de $\mathcal{E}$ estão dentro do retângulo
\begin{figure}[H]
	\centering
	\caption{Retângulo fundamental}
	\includegraphics[scale=0.9]{retangulo-fundamental.pdf}
\end{figure}
Esse retângulo é chamado de \textbf{retângulo fundamental} da elipse $\mathcal{E}$.\index{Elipse!Retângulo Fundamental}

Agora, como $a > b$, então
\[
	\dfrac{x^2}{a^2} + \dfrac{y^2}{a^2} \le \dfrac{x^2}{a^2} + \dfrac{y^2}{b^2} \le \dfrac{x^2}{b^2} + \dfrac{y^2}{b^2}.
\]
Assim, todo ponto $P(x,y)$ de $\mathcal{E}$ satisfaz
\[
	\dfrac{x^2}{a^2} + \dfrac{y^2}{a^2} \le 1 \le \dfrac{x^2}{b^2} + \dfrac{y^2}{b^2}
\]
e com isso $b^2 \le x^2 + y^2 \le a^2$, ou seja, $b \le d(P,O) \le a$. Portanto a elipse está contida entre as circunferências de raio $a$ e $b$. Essa região é chamada de \textbf{coroa fundamental} da elipse.\index{Elipse!Coroa Fundamental}
\begin{figure}[!h]
	\centering
	\caption{Coroa Fundamental}
	\includegraphics[scale=0.9]{coroa-fundamental.pdf}
\end{figure}

Logo os pontos da elipse estão na região compreendida entre o retângulo fundamental e a coroa fundamental.
\begin{figure}[H]
	\centering
	\caption{Região onde se encontram os pontos de uma elipse}
	\includegraphics[scale=0.9]{regiao-fundamental.pdf}
\end{figure}
Assim observe que se $b/a$ estiver próximo de 1, então o retângulo fundamental será próximo de um quadrado e daí a forma da elipse se aproximará de um
circunferência. Por outro lado, quanto mais próximo de 0 estiver $b/a$, mas alongada será a elipse. Mas, sabemos que
\[
	a^2 = b^2 + c^2
\]
daí
\[
	\left(\dfrac{b}{a}\right)^2 + \left(\dfrac{c}{a}\right)^2 = 1.
\]
Dessa equação vemos que se $b/a \to 1$, então $c/a \to 0$ e daí os focos de $\mathcal{E}$ estão ``mais próximos'' de seu centro do que de seus vértices e desse
modo a elipse será mais parecida com um circunferência. Por outro lado, se $b/a \to 0$, então $c/a \to 1$ e assim os focos de $\mathcal{E}$ estão ``mais próximos'' dos vértices do que do centro de $\mathcal{E}$ e desse modo a elipse será mais alongada.  Portanto o quociente $c/a$ mede o quanto os focos de $\mathcal{E}$ estão ``fora do centro''. A essa razão damos o nome de \textbf{excentricidade} (\textit{ex centrum}) da elipse e representamos por\index{Elipse!Excentricidade}
\[
	e = \dfrac{c}{a}.
\]

% subsection elipse (end)

\subsection{Hipérbole} % (fold)
\label{sub:hiperbole}

\begin{definicao}
	Sejam $F_1$ e $F_2$ pontos distintos do plano tais que sua distância seja $2c$. Considere um número real $a$ tal que $0 < a < c$. O lugar geométrico $\mathcal{H}$ dos pontos $P(x,y)$ tais que
	\begin{equation}\label{equacaohiperbole}
		| d(P,F) - d(P,F_2)| = 2a
	\end{equation}
	chama-se \textbf{hipérbole}. Os pontos $F_1$ e $F_2$ são chamados \textbf{focos} da hipérbole, o segmento $F_1F_2$ é chamado \textbf{segmento focal}, o ponto médio do segmento $F_1F_2$ é chamada de \textbf{centro} da hipérbole e o número $2c$ é chamado de \textbf{distância focal}. A reta passando por $F_1$ e $F_2$ é chamada \textbf{reta focal}.\index{Hipérbole} \index{Hipérbole!Focos} \index{Hipérbole!Segmento Focal} \index{Hipérbole!Centro} \index{Hipérbole!Distância Focal} \index{Hipérbole!Reta Focal}
\end{definicao}

Vamos fixar $F_1(-c,0)$ e $F_2(c,0)$. Assim um ponto $P(x,y)$ pertence á hipérbole $\mathcal{H}$ de focos $F_1$ e $F_2$ se
\begin{align*}
	 & | d(P,F_1) - d(P,F_2) | = 2a                                                           \\
	 & d(P,F_1) - d(P,F_2) = \pm 2a                                                           \\
	 & d(P,F_1) = \pm 2a + d(P,F_2)                                                           \\
	 & \left[\sqrt{(x + c)^2 + y^2}\right]^2 = \left[\pm 2a + \sqrt{(x - c)^2 + y^2}\right]^2 \\
	 & (cx - a^2)^2 = \left[\pm a\sqrt{(x - c)^2 + y^2}\right]^2                              \\
	 & (c^2 - a^2)x^2 - a^2y^2 = a^2(c^2 - a^2).
\end{align*}
Fazendo $b^2 = c^2 - a^2$ e dividindo por $a^2b^2$ obtemos
\begin{equation}
	\dfrac{x^2}{a^2} - \dfrac{y^2}{b^2} = 1
\end{equation}
Como no caso da elipse, também temos as relações
\begin{align}
	c^2 = a^2 + b^2                     \\
	c > b > 0\label{pitagorashiperbole} \\
	c > a > 0.
\end{align}
Portanto, se um ponto $P(x,y)$ pertence á hipérbole $\mathcal{H}$, então ele satisfaz a equação \eqref{equacaohiperbole}.

Agora, dado um ponto $A(\alpha,\beta)$ que satisfaz a equação \eqref{equacaohiperbole}, queremos mostrar que $A$ pertence á hipérbole $\mathcal{H}$. Temos
\begin{align}
	\dfrac{\alpha^2}{a^2} - \dfrac{\beta^2}{b^2} = 1\nonumber \\
	\beta^2 = \dfrac{b^2}{a^2}\alpha^2 - b^2.\label{equacaoauxiliarhiperbole}
\end{align}
Usando \eqref{pitagorashiperbole} e \eqref{equacaoauxiliarhiperbole}:
\begin{align*}
	d^2(P,F_1) & = (\alpha + c)^2 + \beta^2                                    \\
	           & = \alpha^2 + 2\alpha c + c^2 + \dfrac{b^2}{a^2}\alpha^2 - b^2
\end{align*}
donde
\[
	d^2(A,F_1) = \left(\dfrac{c}{a}\alpha + a\right)^2.
\]
Do mesmo modo, obtemos
\[
	d^2(A,F_2) = \left(\dfrac{c}{a}\alpha - a\right)^2.
\]
Assim
\begin{align}
	d(A,F_1) = \left|\dfrac{c}{a}\alpha + a\right| \\
	d(A,F_2) = \left|\dfrac{c}{a}\alpha - a\right|.
\end{align}

Mas
\[
	\dfrac{\alpha^2}{a^2} - \dfrac{\beta^2}{b^2} = 1
\]
assim
\begin{equation}\label{pontohiperbole}
	\alpha^2/a^2 \ge 1,
\end{equation}
isto é, $\alpha \le -a$ ou $\alpha \ge a$. Temos então dois casos para analisar:
\begin{enumerate}
	\item Se $\alpha \le -a$, então
	      \[
		      \dfrac{c}{a}\alpha \le -c < -a.
	      \]
	      Portanto, $(c/a)\alpha + a < 0$ e $(c/a)\alpha - a < 0$, donde
	      \begin{align*}
		      d(A,F_1) = -\dfrac{c}{a}\alpha - a \\
		      d(A,F_2) = -\dfrac{c}{a}\alpha + a.
	      \end{align*}
	      e então
	      \[
		      | d(P,F_1) - d(P,F_2)| = |-2a| = 2a.
	      \]
	\item Se $\alpha \ge -a$, então
	      \[
		      \dfrac{c}{a}\alpha \le c > a.
	      \]
	      Portanto, $(c/a)\alpha + a > 0$ e $(c/a)\alpha - a > 0$, donde
	      \begin{align*}
		      d(A,F_1) = \dfrac{c}{a}\alpha + a \\
		      d(A,F_2) = \dfrac{c}{a}\alpha - a.
	      \end{align*}
	      e então
	      \[
		      | d(P,F_1) - d(P,F_2)| = |2a| = 2a.
	      \]
\end{enumerate}
Portanto, se $A(\alpha,\beta)$ satisfaz a equação \eqref{equacaohiperbole}, então $A$ pertence á hipérbole de focos $F_1(-c,0)$ e $F_2(c,0)$. Assim provamos a seguinte proposição:

\begin{proposicao}
	Um ponto $P(x, y)$ é um ponto da hipérbole
	\[
		\mathcal{H}:\ \dfrac{x^2}{a^2} - \dfrac{y^2}{b^2} = 1
	\]
	se, e somente se, as distâncias de $P$ aos focos são:
	\begin{align*}
		d(A,F_1) = \left|\dfrac{c}{a}\alpha + a\right| \\
		d(A,F_2) = \left|\dfrac{c}{a}\alpha - a\right|.
	\end{align*}
\end{proposicao}

Os números $a$, $b$ e $c$ são chamados \textbf{parâmetros geométricos} da hipérbole $\mathcal{H}$ e a equação
\[
	\dfrac{x^2}{a^2} - \dfrac{y^2}{b^2} = 1
\]
é chamada de \textbf{equação reduzida} da hipérbole.\index{Hipérbole!Parâmetros Geométricos} \index{Hipérbole!Equação Reduzida}

\begin{observacao}
	Seja $\mathcal{H} : \dfrac{x^2}{a^2} - \dfrac{y^2}{b^2} = 1$ uma hipérbole de focos no eixo $x$.
	\begin{enumerate}
		\item Nenhum ponto $P(x,y)$ de $\mathcal{H}$ é tal que $-a < x < a$, pois caso contrário teríamos $x^2 < a^2$ o que contradiz a inequação \eqref{pontohiperbole}. Para a ordenada $y$ não há restrição, isto é, qualquer que seja $y$ existe $x$ tal que $\dfrac{x^2}{a^2} - \dfrac{y^2}{b^2} = 1$. Assim a hipérbole não é limitada.
		\item Se $(x,y)$ pertence a $H$, então $(-x,-y)$, $(x,-y)$ e $(-x,y)$ também pertencem a $\mathcal{H}$. Logo a hipérbole é simétrica em relação ao eixo
		      $x$, ao eixo $y$  e em relação à origem do plano cartesiano.
		\item O ponto $(x,0)$ pertence à hipérbole $\mathcal{H}$ se, e só se, $x = \pm a$. Por isso a interseção de $\mathcal{H}$ com a reta focal são os pontos $A_1(-a,0)$ e $A_2(a,0)$, chamados de \textbf{vértices} de $\mathcal{H}$.\index{Hipérbole!Vértices}
	\end{enumerate}
\end{observacao}

Considere a hipérbole $\mathcal{H}$ de equação reduzida
\[
	\mathcal{H} : \dfrac{x^2}{a^2} - \dfrac{y^2}{b^2} = 1
\]
com focos $F_1(-c,0)$ e $F_2(c,0)$. Dada uma reta $r$ de equação $y = mx$, quais as condições sobre $a$, $b$, $c$ e $m$ para que a reta $r$ intercepte $\mathcal{H}$?

Seja $P(x_0,y_0)$ um ponto de interseção entre $r$ e $\mathcal{H}$. Então temos
\begin{align*}
	\dfrac{x_0^2}{a^2} - \dfrac{y_0^2}{b^2} = 1 \\
	y_0 = mx_0,
\end{align*}
isto é,
\[
	\dfrac{x_0^2}{a^2} - \dfrac{m^2x_0^2}{b^2} = 1.
\]
Resolvendo em $x_0$ obtemos
\begin{equation}\label{intersecaohiperbolereta}
	x_0 = \pm \dfrac{ab}{\sqrt{b^2 - a^2m^2}} = \pm \dfrac{b}{\sqrt{\dfrac{b^2}{a^2} - m^2}}.
\end{equation}
Assim
\[
	b^2 - a^2m^2 > 0 \Leftrightarrow m^2 < \dfrac{b^2}{a^2}.
\]
Como $a$ e $b$ são positivos
\[
	-\dfrac{b}{a} < m < \dfrac{b}{a}.
\]
Logo a reta $y = mx$ intercepta a hipérbole $\mathcal{H}$ se, e só se, seu coeficiente angular está entre $-b/a$ e $b/a$. Portanto as retas
\[
	y = -\dfrac{b}{a}x \quad \mbox{e}\quad y = \dfrac{b}{a}x
\]
não interceptam a hipérbole $\mathcal{H}$. Tais retas são chamadas de \textbf{assíntotas da hipérbole}. Mais ainda, da equação \eqref{intersecaohiperbolereta} quando $m \to b/a$, o ponto $x_0$ tende para mais ou menos infinito. Logo os pontos da hipérbole se aproximam das assíntotas a medida que $x$ se afasta da origem.\index{Hipérbole!Assíntotas}

Como no caso da elipse, utilizando a simetria e as assíntotas da hipérbole, podemos mostrar a partir da equação
\[
	y = \dfrac{b}{a}\sqrt{x^2 - a^2}, \quad x \ge a
\]
que a forma da hipérbole será dada pela Figura \ref{FormageralHiperbole}.
\begin{figure}[!h]
	\centering
	\caption{Hipérbole $\dfrac{x^2}{a^2} - \dfrac{y^2}{b^2} = 1$  com assíntotas $y = -\dfrac{b}{a}x$ e $y = \dfrac{b}{a}x$}
	\label{FormageralHiperbole}
	\includegraphics{hiperbole-focos-eixo-X.pdf}
\end{figure}

\begin{proposicao}
	Um equação da forma
	\begin{equation}\label{hiperbolegeral}
		\dfrac{x^2}{p} + \dfrac{y^2}{q} = 1
	\end{equation}
	descreve uma hipérbole se, e somente se, os números reais $p$ e $q$ são de sinal contrário.
\end{proposicao}

\begin{corolario}
	Sejam $p$ e $q$ números reais de sinal contrário e $\mathcal{H}$ a hipérbole de equação \eqref{hiperbolegeral} e parâmetros geométricos $a$ e $b$.
	\begin{enumerate}
		\item Se $p > 0$ e $q < 0$, então $a^2 = p$ e $b^2 = -q$, $\mathcal{H}$ tem centro na origem, focos no eixo $x$ e assíntotas dadas por $y = \pm(b/a)x$.
		\item Se $p < 0$ e $q > 0$, então $a^2 = q$ e $b^2 = -p$, $\mathcal{H}$ tem centro na origem, focos no eixo $y$ e assíntotas dadas por $x = \pm(a/b)y$.
	\end{enumerate}
\end{corolario}

\begin{exemplos}
	\begin{enumerate}
		\item Encontre a equação da hipérbole cujos vértices são $(-15,0)$ e $(15,0)$ e as assíntotas têm equações $5y - 4x = 0$ e $5y + 4x = 0$.
		      \begin{solucao}
			      As assíntotas são dadas por
			      \[
				      y = -\dfrac{b}{a}x \quad\mbox{e }\quad y = \dfrac{b}{a}x.
			      \]
			      A partir dos vértices encontramos $a = 15$. Daí $b = 12$ e então a equação da hipérbole é
			      \[
				      \dfrac{x^2}{225} - \dfrac{y^2}{144} = 1.
			      \]
		      \end{solucao}
		\item Encontre a equação da hipérbole de focos $F_1(0,-\sqrt{20})$ e $F_1(0,\sqrt{20})$ e tal que $a = 2$.
		      \begin{solucao}
			      Como os focos não estão no eixo $x$, vamos usar a equação \eqref{equacaohiperbole}:
			      \begin{align*}
				       & \mid d(P,F) - d(P,F_2)\mid = 4                                                                        \\
				       & \left[\sqrt{x^2 + (y + \sqrt{20})^2}\right]^2 = \left[\pm 4 + \sqrt{x^2 + (y - \sqrt{20})^2}\right]^2 \\
				       & (y\sqrt{20} - 4)^2 = \left[\pm 2\sqrt{x^2 + (y - \sqrt{20})^2}\right]^2                               \\
				       & 20y^2 - 8\sqrt{20}y + 16 = 4x^2 + 4y^2 - 8\sqrt{20}y + 80                                             \\
				       & 16y^2 - 4x^2 = 64.
			      \end{align*}
			      Logo a hipérbole procurada tem equação
			      \[
				      \dfrac{y^2}{4} - \dfrac{x^2}{16} = 1.
			      \]
			      Portanto as assíntotas são
			      \[
				      y = -2x \quad\mbox{e }\quad y = 2x.
			      \]
		      \end{solucao}
	\end{enumerate}
\end{exemplos}

% \begin{observacao}
%   No caso de uma hipérbole $\mathcal{H}$ com focos no eixo $y$, isto é, com equação
%   \[
%     \mathcal{H}: \dfrac{y^2}{a^2} - \dfrac{x^2}{b^2} = 1
%   \]
%   as assíntotas terão equações dadas por
%   \[
%     x = \pm \dfrac{a}{b}y.
%   \]
% \end{observacao}

\begin{figure}[h]
	\centering
	\caption{Hipérbole $\dfrac{y^2}{b^2} - \dfrac{x^2}{a^2} = 1$ com assíntotas $x = -\dfrac{a}{b}y$ e $x = \dfrac{a}{b}y$}
	\includegraphics{hiperbole-focos-eixo-Y.pdf}
\end{figure}

% subsection hiperbole (end)



\subsection{Parábola} % (fold)
\label{sub:parabolas}

\begin{definicao}
	Seja $r$ uma reta e $F$ um ponto que não pertence a $r$. O lugar geométrico $\mathcal{P}$ dos pontos equidistantes de $F$ e $r$ chama-se \textbf{parábola}. O ponto $F$ é chamado de \textbf{foco da parábola} e $r$ de \textbf{reta diretriz}.\index{Parábola} \index{Parábola!Focos} A reta contendo o foco $F$ e perpendicular á reta diretriz é chamada de \textbf{reta focal}.\index{Parábola!Reta diretriz}\index{Parábola!Reta focal}
\end{definicao}


Vamos fixar $F(p,0)$ e $r: x = -p$, onde $p > 0$. Assim um ponto $A(x,y)$ pertence á parábola de foco $F$ e reta diretriz $r$ se, e somente se,
\begin{figure}[!h]%parábola
	\centering
	\caption{Definição da Parábola}
	\includegraphics{parabola.pdf}
\end{figure}

\[
	d(A,r) = d(A,F).
\]
Mas
\begin{align*}
	d(A,r) = |x + p| \\
	d(A,F) = \sqrt{(x - p)^2 + y^2}.
\end{align*}
Logo $A$ pertence á parábola $\mathcal{P}$ se, e somente se,
\begin{align*}
	(|x + p|)^2 = (x - p)^2 + y^2 \\
	x^2 + 2px + p^2 = x^2 - 2px + p^2 + y^2.
\end{align*}
Portanto $A(x,y)$ pertence á parábola $\mathcal{P}$ se, e somente se,
\begin{equation}\label{equacaoparabola}
	y^2 = 4px
\end{equation}

A equação \eqref{equacaoparabola} é chamada de \textbf{equação reduzida} da parábola $\mathcal{P}$. Indica-se\index{Parábola!Equação reduzida}
\[
	\mathcal{P}:\ y^2 = 4px.
\]

\begin{observacao}
	\begin{enumerate}
		\item Se $(x,y)$ satisfaz a equação \eqref{equacaoparabola}, então $x \ge 0$, isto é, nenhum ponto de $\mathcal{P}$ tem abscissa negativa. Já para a ordenada $y$, não existem restrições. Logo a parábola não é limitada.
		\item Se $(x,y)$ pertence a $\mathcal{P}$, então $(x,-y)$ também pertence a $\mathcal{P}$. Logo a parábola é simétrica em relação ao eixo $x$.
		\item O único ponto de interseção de $\mathcal{P}$ com os eixos coordenados é o ponto $(0,0)$. Tal ponto é chamado de \textbf{vértice} da parábola.\index{Parábola!Vértice}
	\end{enumerate}
\end{observacao}

Na equação \eqref{equacaoparabola} podemos isolar $x$ e escrever
\[
	x = \dfrac{y^2}{4p}
\]
obtendo assim $x$ como uma função de $y$. Logo o gráfico da parábola $\mathcal{P}$ é:

\begin{figure}[!h]%parábola
	\centering
	\caption{Parábola $\mathcal{P}: y^2 = 4px$ e reta diretriz $r: x = -p$, $p > 0$}
	\includegraphics{parabola-foco-eixo-X-positivo.pdf}
\end{figure}

Agora, se a diretriz de $\mathcal{P}$ tem equação $r: x = p$ e o foco é o ponto $F(-p,0)$, com $p > 0$, obtemos a equação
\[
	y^2 = -4px.
\]
Neste caso, o gráfico da parábola é dado pela Figura \ref{FormageralParabola}.
\begin{figure}[!h]%parábola
	\centering
	\caption{Parábola $\mathcal{P}: y^2 = -4px$ e reta diretriz $r: x = p$, $p > 0$}
	\label{FormageralParabola}
	\includegraphics{parabola-foco-eixo-X-negativo.pdf}
\end{figure}

Nos demais casos temos:
\begin{figure}[!h]%parábola
	\centering
	\caption{Parábola $\mathcal{P}: x^2 = -4py$ e reta diretriz $r: y = p$, $p > 0$}
	\includegraphics{parabola-foco-eixo-Y-negativo.pdf}
\end{figure}

\begin{figure}[!h]%parábola
	\centering
	\caption{Parábola $\mathcal{P}: x^2 = 4py$ e reta diretriz $r: y = -p$, $p > 0$}
	\includegraphics{parabola-foco-eixo-Y-positivo.pdf}
\end{figure}

\begin{proposicao}
	As equações $y^2 = qx$ e $x^2 = qy$ descrevem uma parábola se, e somente se, $q \ne 0$.
\end{proposicao}
% subsection parabolas (end)
% section conicas (end)

\section{Rotação e Translação de Eixos} % (fold)
\label{sec:rotacao_e_translacao_de_eixos}


Para determinar a cônica representada pela equação
\[
	g(x, y) = ax^2 + bxy + cy^2 + dx + ey + f = 0
\]
iremos utilizar translações e rotações de eixos para simplificar sua equação.

\subsection{Translação de eixos} % (fold)
\label{sub:translacaoo_de_eixos}
Considere a seguinte cônica
\begin{equation}\label{conicaexemplo}
	x^2 + 2y^2 - 4x - 4y - 1 = 0.
\end{equation}
Podemos escrever
\begin{align*}
	 & (x^2 - 4x + 4) - 4 + 2(y^2 - 2y + 1 - 1) - 1 = 0 \\
	 & (x - 2)^2 - 4 + 2(y - 1)^2 - 2 - 1 = 0           \\
	 & (x - 2)^2 + 2(y - 1)^2 = 7.
\end{align*}
Fazendo a mudança
\begin{align}
	x - 2 = x_1\label{mudancaX} \\
	y - 1 = y_1\label{mudancaY}
\end{align}
obtemos
\[
	\dfrac{x_1^2}{7} + \dfrac{y_1^2}{\dfrac{7}{2}} = 1
\]
que representa uma elipse com focos no eixo $x_1$.

O efeito das equações \eqref{mudancaX} e \eqref{mudancaY} foi o de mudar o centro da elipse de equação \eqref{conicaexemplo} que estava no ponto $(2,1)$ no sistema original para o ponto $(0,0)$ considerando os eixos coordenados $x_1$ e $y_1$.

Utilizando-se as equações \eqref{mudancaX} e \eqref{mudancaY} podemos facilmente escrever as coordenadas de um ponto $P$ qualquer, tanto em relação ao sistema de eixos coordenados $x$ e $y$, quanto ao sistema de eixos coordenados $x_1$ e $y_1$. Por exemplo, o ponto $P$ de coordenadas $P = (4,0)$, em relação ao sistema $xy$, terá coordenadas $(2,-1)$ em relação ao sistema de eixos $x_1y_1$. As mudanças introduzidas pelas equações \eqref{mudancaX} e \eqref{mudancaY} são chamadas de \textbf{translações de eixos}. De modo geral, se $(x,y)$ são as coordenadas de um ponto $P$ em relação ao ponto $(0,0)$, então as coordenadas de $P$ em relação aos eixos $x_1$ e $y_1$, centrados no ponto $(\alpha,\beta)$ são dadas por\index{Translação de Eixos}
\begin{align}
	x_1 = x - \alpha\label{mudancaXgeral} \\
	y_1 = y - \beta\label{mudancaYgeral}.
\end{align}

\begin{figure}[!h]
	\centering
	\caption{Translação de eixos}
	\includegraphics{translacao-eixos.pdf}
\end{figure}
As mudanças introduzidas pelas equações \eqref{mudancaXgeral} e \eqref{mudancaYgeral} permitem remover os termos lineares da equação \eqref{equacaoconica}.

\begin{exemplos}
	Identifique as cônicas:
	\begin{enumerate}
		\item $4x^2 - 8x + 9y^2 - 36y + 4 = 0$
		      \begin{solucao}
			      Completando quadrados
			      \begin{align*}
				       & 4(x^2 - 2x + 1 - 1) + 9(y^2 - 4y + 4 - 4) + 4 = 0 \\
				       & 4(x - 1)^2 + 9(y - 2)^2 = 36.
			      \end{align*}
			      Fazendo
			      \begin{align*}
				      x - 1 = x_1 \\
				      y - 2 = y_1
			      \end{align*}
			      obtemos
			      \[
				      \dfrac{x_1^2}{9} + \dfrac{y_1^2}{4} = 1
			      \]
			      que representa uma elipse com focos no eixo $x_1$. No novo sistema de coordenadas centrado no ponto os vértices são
			      \begin{align*}
				      \overline{A_1}(-3,0),\quad \overline{A_2}(3,0) \\
				      \overline{B_1}(0,-2),\quad \overline{B_2}(0,2).
			      \end{align*}
			      No sistema original os vértices são
			      \begin{align*}
				      A_1(-2,2),\quad A_2(2,2) \\
				      B_1(1,0),\quad B_2(1,4).
			      \end{align*}
			      \begin{figure}[!h]
				      \centering
				      \caption{Elipse $4x^2 - 8x + 9y^2 - 36y + 4 = 0$}
				      \includegraphics{elipse-eixos-tranladados.pdf}
			      \end{figure}
		      \end{solucao}
		\item $x^2 - 2y^2 - 6x - 8y - 1 = 0$
		      \begin{solucao}
			      Completando quadrados
			      \begin{align*}
				       & (x^2 - 6x + 9 - 9) - 2(y^2 + 4y + 4 - 4) - 1 = 0 \\
				       & (x - 3)^2 - 2(y + 2)^2 = 2.
			      \end{align*}
			      Fazendo
			      \begin{align*}
				      x - 3 = x_1 \\
				      y + 2 = y_1
			      \end{align*}
			      obtemos
			      \[
				      \dfrac{x_1^2}{2} - y_1^2 = 1
			      \]
			      que representa uma hipérbole com focos no eixo $x_1$. Suas assíntotas são
			      \[
				      y_1 = \pm \dfrac{x_1}{\sqrt{2}}.
			      \]
			      Os vértices são $\overline{A_1}(-\sqrt{2},0)$ e $\overline{A_2}(\sqrt{2},0)$.
			      No sistema original suas assíntotas são
			      \begin{align*}
				      y = \dfrac{x}{\sqrt{2}} - \left(\dfrac{3}{\sqrt{2}} + 2\right) \\
				      y = -\dfrac{x}{\sqrt{2}} + \left(\dfrac{3}{\sqrt{2}} - 2\right).
			      \end{align*}
			      Os vértices são $A_1(-3 - \sqrt{2},-2)$ e $A_2(-3 + \sqrt{2},-2)$.
			      \begin{figure}[!h]
				      \centering
				      \caption{Hipérbole $x^2 - 2y^2 - 6x - 8y - 1 = 0$}
				      \includegraphics{hiperbole-eixos-tranladados.pdf}
			      \end{figure}
		      \end{solucao}
		\item $x^2 - y^2 + 2x - 6y - 8 = 0$
		      \begin{solucao}
			      Completando quadrados
			      \begin{align*}
				       & (x^2 - 2x + 1 - 1) - (y^2 + 6y + 9 - 9) - 8 = 0 \\
				       & (x - 1)^2 - (y + 3)^2 = 0                       \\
				       & [(x - 1) - (y + 3)][(x - 1) + (y + 3)] = 0
			      \end{align*}
			      que representam as retas
			      \begin{align*}
				      r: x - y - 4 = 0 \\
				      s: x + y + 2 = 0
			      \end{align*}
			      que se interceptam no ponto $(1,-3)$.
		      \end{solucao}
	\end{enumerate}
\end{exemplos}
% % subsection translacaoo_de_eixos (end)

\subsection{Rotação de eixos} % (fold)
\label{sub:rotacao_de_eixos}

Considere o sistema cartesiano com centro em $(0,0)$. Sejam $\vec{e_1} = (1,0)$ e $\vec{e_2} = (0,1)$. Para qualquer vetor $\vec{P} = (x,y)$ podemos escrever
\begin{align}\label{equacaoxy}
	\vec{P} = (x,y) = (x,0) + (0,y) = x(1,0) + y(0,1) = x\vec{e_1} + y\vec{e_2}.
\end{align}

\begin{figure}[!h]
	\centering
	\caption{Rotação de eixos: vetores $\vec{e_1}$ e $\vec{e_2}$}
	\includegraphics{rotacao-eixos-vetores-base.pdf}
\end{figure}

Agora, efetuando-se uma rotação no sentido anti-horário nos eixos $x$ e $y$ de um ângulo $\theta$, obtemos novos eixos coordenados $x_1$ e $y_1$. Tome vetores unitários $\vec{u_1}$ e $\vec{u_1}$ nos eixos $x_1$ e $y_1$. Em relação aos eixos originais podemos escrever
\begin{figure}[!h]
	\centering
	\caption{Rotação dos eixos por um ângulo $\theta$}
	\includegraphics{rotacao-eixos.pdf}
\end{figure}
\begin{align*}
	 & \vec{u_1} = (\cos\theta,\sin\theta) = \cos\theta\vec{e_1} + \sin\theta\vec{e_2}    \\
	 & \vec{u_2} = (-\sin\theta,\cos\theta) = -\sin\theta\vec{e_1} + \cos\theta\vec{e_2}.
\end{align*}
Mas, os vetores $\vec{u_1}$ e $\vec{u_2}$ são unitários e ortogonais, assim podemos escrever o vetor $\vec{P}$ como combinação de escalares nos eixos rotacionados, isto é,
\begin{equation}\label{equacaox1y1}
	\vec{P} = x_1\vec{u_1} + y_1\vec{u_2}.
\end{equation}

Daí de \eqref{equacaoxy} e \eqref{equacaox1y1} obtemos
\[
	x\vec{e_1} + y\vec{e_2} = x_1\vec{u_1} + y_1\vec{u_2}
\]
e substituindo $\vec{u_1}$ e $\vec{u_2}$
\[
	x\vec{e_1} + y\vec{e_2} = (x_1\cos\theta - y_1\sin\theta)\vec{e_1} + (x_1\sin\theta + y_1\cos\theta)\vec{e_2}.
\]
Logo
\begin{align}
	 & x = x_1\cos\theta - y_1\sin\theta\label{rotacaoX} \\
	 & y = x_1\sin\theta + y_1\cos\theta\label{rotacaoY}
\end{align}
e então isolando $x_1$ e $y_1$
\begin{align*}
	 & x_1 = x\cos\theta + y\sin\theta   \\
	 & y_1 = -x\sin\theta + y\cos\theta.
\end{align*}
% Portanto para eliminar o termo $xy$ da equação $g(x,y) = 0$, substituímos \eqref{rotacaoX} e \eqref{rotacaoY} em $g(x,y) = 0$ e determinanos o ângulo $\theta$ que elimina o termo $xy$.

\begin{exemplos}
	Seja $P$ o ponto $P(6,4)$. Efetuando-se uma rotação de um ângulo de $\theta = \pi/3$ radianos nos eixos, as coordenadas de $P$ em relação aos novos eixos são
	\begin{align*}
		x_1 = 6\cos(\pi/3) + 4\sin(\pi/3) = 3 + 2\sqrt{3} \\
		y_1 = -6\sin(\pi/3) + 4\cos(\pi/3) = 2 - 3\sqrt{3}.
	\end{align*}
	Daí no novo sistema $P(3 + 2\sqrt{3}, 2 - 3\sqrt{3})$.
\end{exemplos}

% \begin{exemplos}
%   Identifique as cônicas:
%   \begin{enumerate}
%     \item $3x^2 + 3y^2 - 10xy + 12\sqrt{2}x - 4\sqrt{2}y + 32 = 0$
%     \begin{solucao}
%       Fazendo as mudanças
%       \begin{align}
%         x = x_1\cos\theta - y_1\sin\theta\\
%         x = x_1\sin\theta + y_1\cos\theta
%       \end{align}
%       temos
%       \begin{align*}
%         &3(x_1\cos\theta - y_1\sin\theta)^2 + 3(x_1\sin\theta + y_1\cos\theta)^2\\ &- 10(x_1\cos\theta - y_1\sin\theta)(x_1\sin\theta + y_1\cos\theta)\\ & + 12\sqrt{2}(x_1\cos\theta - y_1\sin\theta) - 4\sqrt{2}(x_1\sin\theta + y_1\cos\theta)\\ &= (3\cos^2\theta + 3\sin^2\theta - 10\sin\theta\cos\theta)x_1^2 \\ &+ (3\sin^2\theta + 3\cos^2\theta + 10\sin\theta\cos\theta)y_1^2 \\ &+ (6\sin\theta\cos\theta - 10\cos^2\theta + 10\sin^2\theta - 6\sin\theta\cos\theta)x_1y_1 \\ &+ (12\sqrt{2}\cos\theta - 4\sqrt{2}\sin\theta)x_1 + (-12\sqrt{2}\sin\theta - 4\sqrt{2}\cos\theta)y_1 + 32 = 0.
%       \end{align*}
%       Assim $\theta$ deve ser tal que
%       \[
%         -10\cos^2\theta + 10\sin^2\theta = 0,
%       \]
%       isto é, $\theta = \pi/4$. Substituindo esse valor de $\theta$ na equação anterior e simplificando obtemos
%       \[
%         x_1^2 - 4y_1^2 - 4x_1 + 8y_1 - 16 = 0.
%       \]
%       Nessa nova equação podemos completar quadrados
%       \begin{align*}
%         &(x_1^2 - 4x_1 + 4) - 4 - 4(y_1^2 - 2y_1 + 1 - 1) - 16 = 0\\
%         &(x_1 - 2)^2 - 4(y_1 - 1)^2 = 16.
%       \end{align*}
%       Fazendo
%       \begin{align*}
%         x_2 = x_1 - 2\\
%         y_2 = y_1 - 1
%       \end{align*}
%       encontramos
%       \[
%         \dfrac{x_2^2}{16} - \dfrac{y_2^2}{4} = 1
%       \]
%       que é uma hipérbole de vértices no eixo $x_2$.
%       \begin{figure}[!h]
%         \centering
%         \caption{Hipérbole $3x^2 +3y^2 - 10xy + 12\sqrt{2}x - 4\sqrt{2}y + 32 = 0$}
%         \includegraphics{hiperbole-rotacionada-transladada.pdf}
%       \end{figure}
%     \end{solucao}
%     \item $52x^2 - 72xy + 73y^2 - 400 = 0$
%     \begin{solucao}
%       Fazendo as mudanças
%       \begin{align}
%         x = x_1\cos\theta - y_1\sin\theta\\
%         x = x_1\sin\theta + y_1\cos\theta
%       \end{align}
%       temos
%       \begin{align}\label{conicaelipse}
%         &52(x_1\cos\theta - y_1\sin\theta)^2 + 73(x_1\sin\theta + y_1\cos\theta)^2\nonumber \\ &- 72(x_1\cos\theta - y_1\sin\theta)(x_1\sin\theta + y_1\cos\theta) - 400 = 0\nonumber\\ &(52\cos^2\theta - 72\sin\theta\cos\theta + 73\sin^2\theta)x_1^2 + (52\sin^2\theta + 72\sin\theta\cos\theta + 73\cos^2\theta)y_1^2 \nonumber\\&+ (-104\sin\theta\cos\theta + 72\sin^2\theta - 72\cos^2\theta + 146\sin\theta\cos\theta)x_1y_1 - 400 = 0
%       \end{align}
%       Assim devemos ter
%       \begin{align*}
%         &-104\sin\theta\cos\theta + 72\sin^2\theta - 72\cos^2\theta + 146\sin\theta\cos\theta = 0\\
%         &42\sin\theta\cos\theta + 72(\sin^2\theta - \cos^2\theta) = 0\\
%         &21\sin(2\theta) - 72\cos(2\theta) = 0\\
%         &\tan(2\theta) = \dfrac{24}{7}.
%       \end{align*}
%       Agora
%       \[
%         \cos(2\theta) = \dfrac{1}{\sqrt{1 + \tan^2(2\theta)}}
%       \]
%       e usando as equações $\cos^2\theta = (1/2)(1 + \cos(2\theta))$ e $\sin^2\theta = (1/2)(1 - \cos(2\theta))$ obtemos
%       \[
%         \cos\theta = \dfrac{4}{5} \quad \mbox{e}\quad \sin\theta = \dfrac{3}{5}.
%       \]
%       Substituindo na equação \eqref{conicaelipse} obtemos
%       \begin{align*}
%         &\left(52\dfrac{16}{25} - 72\dfrac{12}{25} + 73\dfrac{9}{25}\right)x_1^2 + \left(52\dfrac{9}{25} + 72\dfrac{12}{25} + 73\dfrac{16}{25}\right)y_1^2 - 400 = 0\\
%         &25x_1^2 + 100y_1^2 = 400\\
%         &\dfrac{x_1^2}{16} + \dfrac{y_1^2}{4} = 1
%       \end{align*}
%       que representa uma elipse com focos no eixo $x_1$.
%       \begin{figure}[!h]
%         \centering
%         \caption{Elipse $52x^2 - 72xy + 73y^2 - 400 = 0$}
%         \includegraphics{elipse-rotacionada.pdf}
%       \end{figure}
%     \end{solucao}
%     \item $x^2 - 2xy + y^2 - 2x - 2y + 1 = 0$
%     \begin{solucao}
%       Fazendo as mudanças
%       \begin{align}
%         x = x_1\cos\theta - y_1\sin\theta\\
%         x = x_1\sin\theta + y_1\cos\theta
%       \end{align}
%       temos
%       \begin{align}\label{conicaparabola}
%         &(x_1\cos\theta - y_1\sin\theta)^2 - 2(x_1\cos\theta - y_1\sin\theta)(x_1\sin\theta + y_1\cos\theta)\nonumber \\&+ (x_1\sin\theta + y_1\cos\theta)^2 - 2(x_1\cos\theta - y_1\sin\theta) - 2(x_1\sin\theta + y_1\cos\theta) + 1 = 0\nonumber\\
%         &(\cos^2\theta - 2\sin\theta\cos\theta + \sin^2\theta)x_1^2 + (\sin^2\theta + 2\sin\theta\cos\theta + \cos^2\theta)y_1^2\nonumber\\ &+ (-2\sin\theta\cos\theta - 2\cos^2\theta + 2\sin^2\theta + 2\sin\theta\cos\theta)x_1y_1\nonumber\\ &+ (-2\cos\theta - 2\sin\theta)x_1 + (2\sin\theta - 2\cos\theta)y_1 + 1 = 0.
%       \end{align}
%       Assim devemos ter
%       \[
%         -2\cos^2\theta + 2\sin^2\theta = 0,
%       \]
%       isto é, $\theta = \pi/4$. Substituindo em \eqref{conicaparabola} obtemos a equação
%       \[
%         y_1^2 = \sqrt{2}\left(x_1 - \dfrac{1}{2\sqrt{2}}\right).
%       \]
%       Fazendo
%       \begin{align*}
%         x_1 - \dfrac{1}{2\sqrt{2}} = x_2\\
%         y_1 = y_2
%       \end{align*}
%       encontramos
%       \[
%         y_2^2 = \sqrt{2}x_2
%       \]
%       que representa uma parábola com foco no eixo $x_2$ e reta diretriz $x_2 = -\sqrt{2}/4$.
%       \begin{figure}[!h]
%         \centering
%         \caption{Parábola $x^2 - 2xy + y^2 - 2x - 2y + 1 = 0$}
%         \includegraphics{parabola-rotacionada-transladada.pdf}
%       \end{figure}
%     \end{solucao}
%   \end{enumerate}
% \end{exemplos}

% subsection rotacao_de_eixos (end)

% \section{Rotações de Eixos: método matricial} % (fold)
% \label{sec:rotacoes_de_eixos_metodo_matricial}

Considere uma cônica de equação
\begin{equation}\label{equacao_geral_conica}
	ax^2 + bxy + cx^2 + dx + ey + f = 0.
\end{equation}

Queremos remover o termo $xy$ da equação \eqref{equacao_geral_conica} para então usar uma translação de eixos, se necessário, e com isso identificar a cônica definida por tal equação. Para isso vamos tentar encontrar um ângulo $\theta$ tal que a equação \eqref{equacao_geral_conica} possa ser transformada numa equação da forma
\begin{equation}\label{equacao_rotacionada_conica}
	a'x_1^2 + c'y_1^2 + d'x_1 + e'y_1 + f' = 0.
\end{equation}

Inicialmente observe que a \eqref{equacao_geral_conica} pode ser escrita na forma
\begin{equation}\label{equacao_matricial_conica}
	X^tAX + KX + F = [0]
\end{equation}
onde
\[
	A = \begin{bmatrix}
		a   & b/2 \\
		b/2 & c
	\end{bmatrix}
	, \quad K = \begin{bmatrix}
		d & e
	\end{bmatrix}, \quad X = \begin{bmatrix}
		x \\
		y
	\end{bmatrix}\quad\mbox{e}\quad F = [f]
\]
e $t$ denota a matriz transposta. Queremos então encontrar um ângulo $\theta$ tal que a substituição
\begin{equation}\label{substituicao_rotacao_eixos}
	X = \begin{bmatrix}
		\cos\theta & -\sin\theta \\
		\sin\theta & \cos\theta
	\end{bmatrix}\begin{bmatrix}
		x_1 \\y_1
	\end{bmatrix}
\end{equation}
transforma a equação \eqref{equacao_geral_conica} na equação \eqref{equacao_rotacionada_conica}.

Sejam
\[
	R_\theta = \begin{bmatrix}
		\cos\theta & -\sin\theta \\
		\sin\theta & \cos\theta
	\end{bmatrix}
\]
que é chamada de \textbf{matriz de rotação} e
\[
	X_1 = \begin{bmatrix}x_1\\y_1\end{bmatrix}.
\]
Note que
\[
	R_\theta R_\theta^t = \begin{bmatrix}
		\cos\theta & -\sin\theta \\
		\sin\theta & \cos\theta
	\end{bmatrix}\begin{bmatrix}
		\cos\theta  & \sin\theta \\
		-\sin\theta & \cos\theta
	\end{bmatrix} = \begin{bmatrix}
		1 & 0 \\
		0 & 1
	\end{bmatrix} = R_\theta^t R_\theta,
\]
ou seja, $(R_\theta^t)^{-1} = R_\theta$.

Agora, substituindo \eqref{substituicao_rotacao_eixos} em \eqref{equacao_geral_conica} obtemos
\begin{align*}
	(R_\theta X_1)^t A (R_\theta X_1) + KR_\theta X_1 + F = [0] \\
	X_1^t(R_\theta^t A R_\theta)X_1 + (KR_\theta)X_1 + F = [0].
\end{align*}
Note que
\begin{align*}
	R_\theta^t A R_\theta = \begin{bmatrix}
		                        a'   & b'/2 \\
		                        b'/2 & c'
	                        \end{bmatrix} = B \\
	KR_\theta = \begin{bmatrix}
		            d' & e'
	            \end{bmatrix} = K'.
\end{align*}

Mas $I_2 = R_\theta^t R_\theta$ e daí podemos escrever
\begin{align*}
	B - \lambda I_2 = R_\theta^t A R_\theta - \lambda I_2 = R_\theta^t A R_\theta - R_\theta^t (\lambda I_2) R_\theta = R_\theta^t (A - \lambda I_2)R_\theta
\end{align*}
para todo $\lambda \in \real$. Com isso
\begin{align*}
	\det(B - \lambda I_2) = \det[R_\theta^t(A - \lambda I_2)R_\theta] = \det(R_\theta^t)\det(A - \lambda I_2)\det(R_\theta) = \det(A - \lambda I_2).
\end{align*}

Assim escolhendo $\theta$ tal que $b' = 0$, o que é sempre possível fazer, obtemos
\begin{align*}
	\det(A - \lambda I_2) = \det(B - \lambda I_2) = \det \begin{bmatrix}
		                                                     a' - \lambda & 0            \\
		                                                     0            & c' - \lambda
	                                                     \end{bmatrix} = (\lambda - a')(\lambda - c').
\end{align*}

Logo os coeficientes $a'$ e $c'$ são raízes da equação de segundo grau
\begin{equation}
	p(\lambda) = \det(A - \lambda I_2) = \det \begin{bmatrix}
		a - \lambda & b/2         \\
		b/2         & c - \lambda
	\end{bmatrix}.
\end{equation}

Agora para determinar o ângulo $\theta$ partimos de
\[
	B = R_\theta^t A R_\theta
\]
e multiplicamos á esquerda por $R_\theta$ obtendo
\[
	R_\theta B = AR_\theta.
\]
Daí
\begin{align*}
	AR_\theta = A \begin{bmatrix}
		              \cos\theta & -\sin\theta \\
		              \sin\theta & \cos\theta
	              \end{bmatrix} = \begin{bmatrix}
		                              A \begin{bmatrix}
			  \cos\theta \\ \sin\theta
		  \end{bmatrix} & A \begin{bmatrix}
			                    -\sin\theta \\ \cos\theta
		                    \end{bmatrix}
	                              \end{bmatrix}
\end{align*}
e
\begin{align*}
	R_\theta B = \begin{bmatrix}
		             \cos\theta & -\sin\theta \\
		             \sin\theta & \cos\theta
	             \end{bmatrix} \begin{bmatrix}
		                           a' & 0  \\
		                           0  & c'
	                           \end{bmatrix} = \begin{bmatrix}
		                                           a' \begin{bmatrix}
			   \cos\theta \\ \sin\theta
		   \end{bmatrix} & c' \begin{bmatrix}
			                      -\sin\theta \\ \cos\theta
		                      \end{bmatrix}
	                                           \end{bmatrix}.
\end{align*}

Como $R_\theta B = AR_\theta$ então das duas últimas equação vemos que se tomarmos
\[
	U_1 = \begin{bmatrix}
		\cos\theta \\
		\sin\theta
	\end{bmatrix} = (\cos\theta, \sin\theta)
\]
devemos ter
\[
	AU_1 = a'U_1,
\]
isto é,
\begin{align*}
	AU_1 = a'I_2 U_1       \\
	AU_1 - a'I_2 U_1 = [0] \\
	(A - a'I_2)U_1 = [0].
\end{align*}
Assim $U_1$ é uma solução de norma 1 do sistema linear
\[
	(A - a'I_2)X = [0].
\]
A solução $U_2$ é obtida de $U_1$ trocando a posição das componentes de $U_1$ e depois o sinal da primeira posição.

Portanto a rotação que elemina o termo $xy$ é dada por $R_\theta = [U_1\ U_2]$. Assim temos o seguinte teorema:

\begin{teorema}
	Considere a equação
	\[
		ax^2 + bxy + cy^2 + dx + ey + f = 0
	\]
	com $a$, $b$, $c$, $d$, $e$, $f \in \real$, sendo $a$, $b$ e $c$ não necessariamente nulos. Então mediante uma rotação da forma
	\[
		X = R_\theta X_1
	\]
	onde
	\[
		X = \begin{bmatrix}
			x \\y
		\end{bmatrix}, \quad
		R_\theta = \begin{bmatrix}
			\cos\theta & -\sin\theta \\
			\sin\theta & \cos\theta
		\end{bmatrix}, \quad\mbox{e}\quad
		X_1 = \begin{bmatrix}
			x_1 \\y_1
		\end{bmatrix}
	\]
	a equação da cônica acima pode ser transformada em
	\[
		a'x_1^2 + c'y_1^2 + d'x_1 + e'y_1 + f = 0
	\]
	em que $a'$, $c'$ são raízes de
	\[
		p(\lambda) = \det(A - \lambda I_2) = \det \begin{bmatrix}
			a - \lambda & b/2         \\
			b/2         & c - \lambda
		\end{bmatrix}.
	\]

	Mais ainda,
	\[
		U_1 = \begin{bmatrix}
			\cos\theta \\ \sin\theta
		\end{bmatrix}
	\]
	é uma solução de norma 1 do sistema linear
	\[
		\begin{bmatrix}
			a - a' & b/2    \\
			b/2    & c - a'
		\end{bmatrix} \begin{bmatrix}
			x \\y
		\end{bmatrix} = \begin{bmatrix}
			0 \\0
		\end{bmatrix}.
	\]
\end{teorema}

\begin{exemplos}
	Identifique as cônicas:
	\begin{enumerate}
		\item $52x^2 - 72xy + 73y^2 - 400 = 0$
		      \begin{solucao}
			      Inicialmente podemos escrever tal equação na forma
			      \[
				      X^t A X + KX - 400 = 0
			      \]
			      onde
			      \[
				      X = \begin{bmatrix}
					      x \\y
				      \end{bmatrix}, \quad A = \begin{bmatrix}
					      52  & -36 \\
					      -36 & 73
				      \end{bmatrix},\quad K = \begin{bmatrix}
					      0 & 0
				      \end{bmatrix}.
			      \]
			      Fazendo a mudança $X = R_\theta X_1$ onde
			      \[
				      X_1 = \begin{bmatrix}
					      x_1 \\y_1
				      \end{bmatrix} \quad\mbox{e}\quad R_\theta = \begin{bmatrix}
					      \cos\theta & -\sin\theta \\
					      \sin\theta & \cos\theta
				      \end{bmatrix}
			      \]
			      removeremos o termo $xy$. Para isso precisamos encontrar as raízes de
			      \begin{align*}
				      p(\lambda) & = \det (A - \lambda I_2) = \det \begin{bmatrix}
					                                                   52 - \lambda & -36          \\
					                                                   -36          & 73 - \lambda
				                                                   \end{bmatrix} \\ &= (52 - \lambda)(73 - \lambda) - 36^2 \\ &= \lambda^2 - 125\lambda + 2500.
			      \end{align*}
			      Suas raízes são $\lambda_1 = 25$ e $\lambda_2 = 100$. Tomando $a' = 25$ e $c' = 100$ precisamos resolver o sistema
			      \begin{align*}
				      (A - 25I_2)X = \begin{bmatrix}
					                     0 \\0
				                     \end{bmatrix} \\
				      \begin{bmatrix}
					      27  & -36 \\
					      -36 & 48
				      \end{bmatrix} \begin{bmatrix}
					                    x \\y
				                    \end{bmatrix} = \begin{bmatrix}
					                                    0 \\0
				                                    \end{bmatrix}
			      \end{align*}
			      para encontrar a matriz de rotação. A solução deste sistema é $x = (4/3)y$ e para determinar a matriz precisamos escolher soluções de norma 1. Assim podemos tomar, por exemplo,
			      \begin{align*}
				      U_1 = (\cos\theta, \sin\theta) = (4/5, 3/5) \\
				      U_2 = (-\sin\theta, \cos\theta) = (-3/5, 4/5).
			      \end{align*}
			      Daí o ângulo de rotação é $\theta = \arccos(4/5)$, a matriz de rotação é
			      \[
				      R_\theta = \begin{bmatrix}
					      4/5 & -3/5 \\
					      3/5 & 4/5
				      \end{bmatrix}
			      \]
			      e
			      \[
				      B = \begin{bmatrix}
					      25 & 0   \\
					      0  & 100
				      \end{bmatrix}.
			      \]
			      Com isso a equação original torna-se
			      \begin{align*}
				      25x_1^2 + 100y_1^2 = 400 \\
				      \dfrac{x_1^2}{16} + \dfrac{y_1^2}{4} = 1
			      \end{align*}
			      que é uma elipse com focos no eixo $x_1$.
			      \begin{figure}[!h]
				      \centering
				      \caption{Elipse $52x^2 - 72xy + 73y^2 - 400 = 0$}
				      \includegraphics{elipse-rotacionada.pdf}
			      \end{figure}
		      \end{solucao}

		\item $x^2 - 2xy + y^2 - 2x - 2y + 1 = 0$
		      \begin{solucao}
			      Inicialmente podemos escrever tal equação na forma
			      \[
				      X^t A X + KX + 1 = 0
			      \]
			      onde
			      \[
				      X = \begin{bmatrix}
					      x \\y
				      \end{bmatrix}, \quad A = \begin{bmatrix}
					      1  & -1 \\
					      -1 & 1
				      \end{bmatrix},\quad K = \begin{bmatrix}
					      -2 & -2
				      \end{bmatrix}.
			      \]
			      Fazendo a mudança $X = R_\theta X_1$ onde
			      \[
				      X_1 = \begin{bmatrix}
					      x_1 \\y_1
				      \end{bmatrix} \quad\mbox{e}\quad R_\theta = \begin{bmatrix}
					      \cos\theta & -\sin\theta \\
					      \sin\theta & \cos\theta
				      \end{bmatrix}
			      \]
			      removeremos o termo $xy$. Para isso precisamos encontrar as raízes de
			      \begin{align*}
				      p(\lambda) & = \det (A - \lambda I_2) = \det \begin{bmatrix}
					                                                   1 - \lambda & -1          \\
					                                                   -1          & 1 - \lambda
				                                                   \end{bmatrix} \\ &= (1 - \lambda)^2 - 1 = \lambda^2 - 2\lambda.
			      \end{align*}
			      Suas raízes são $\lambda_1 = 0$ e $\lambda_2 = 2$. Tomando $a' = 0$ e $c' = 2$ precisamos resolver o sistema
			      \begin{align*}
				      (A - 0I_2)X = \begin{bmatrix}
					                    0 \\0
				                    \end{bmatrix} \\
				      \begin{bmatrix}
					      1  & -1 \\
					      -1 & 1
				      \end{bmatrix} \begin{bmatrix}
					                    x \\y
				                    \end{bmatrix} = \begin{bmatrix}
					                                    0 \\0
				                                    \end{bmatrix}
			      \end{align*}
			      para encontrar a matriz de rotação. A solução deste sistema é $x = y$ e para determinar a matriz precisamos escolher soluções de norma 1. Assim podemos tomar, por exemplo,
			      \begin{align*}
				      U_1 = (\cos\theta, \sin\theta) = (1/\sqrt{2}, 1/\sqrt{2}) \\
				      U_2 = (-\sin\theta, \cos\theta) = (-1/\sqrt{2}, 1/\sqrt{2}).
			      \end{align*}
			      Daí o ângulo de rotação é $\theta = \arccos(1/\sqrt{2}) = \pi/4$, a matriz de rotação é
			      \[
				      R_\theta = \begin{bmatrix}
					      1/\sqrt{2} & -1/\sqrt{2} \\
					      1/\sqrt{2} & 1/\sqrt{2}
				      \end{bmatrix}
			      \]
			      e
			      \[
				      B = \begin{bmatrix}
					      0 & 0 \\
					      0 & 2
				      \end{bmatrix}.
			      \]
			      Com isso a equação original torna-se
			      \begin{align*}
				      2y_1^2 - 4\sqrt{2}x_1 + 1 = 0 \\
				      y_1^2 = \sqrt{2}\left(x_1 - \dfrac{1}{2\sqrt{2}}\right).
			      \end{align*}
			      Fazendo a mudança
			      \begin{align*}
				      x_2 = x_1 - \dfrac{1}{2\sqrt{2}} \\
				      y_2 = y_1
			      \end{align*}
			      obtemos a equação
			      \[
				      y_2^2 = \sqrt{2}x_2
			      \]
			      que é uma parábola com foco no eixo $x_2$ e reta diretriz $r: x_2 = -\sqrt{2}/4$.
			      \begin{figure}[!h]
				      \centering
				      \caption{Parábola $x^2 - 2xy + y^2 - 2x - 2y + 1 = 0$}
				      \includegraphics{parabola-rotacionada-transladada.pdf}
			      \end{figure}
		      \end{solucao}

		\item $3x^2 - 10xy + 3y^2 - 12\sqrt{2}x - 4\sqrt{2}y + 32 = 0$
		      \begin{solucao}
			      Inicialmente podemos escrever tal equação na forma
			      \[
				      X^t A X + KX + 1 = 0
			      \]
			      onde
			      \[
				      X = \begin{bmatrix}
					      x \\y
				      \end{bmatrix}, \quad A = \begin{bmatrix}
					      3  & -5 \\
					      -5 & 3
				      \end{bmatrix},\quad K = \begin{bmatrix}
					      12\sqrt{2} & -4\sqrt{2}
				      \end{bmatrix}.
			      \]
			      Fazendo a mudança $X = R_\theta X_1$ onde
			      \[
				      X_1 = \begin{bmatrix}
					      x_1 \\y_1
				      \end{bmatrix} \quad\mbox{e}\quad R_\theta = \begin{bmatrix}
					      \cos\theta & -\sin\theta \\
					      \sin\theta & \cos\theta
				      \end{bmatrix}
			      \]
			      removeremos o termo $xy$. Para isso precisamos encontrar as raízes de
			      \begin{align*}
				      p(\lambda) & = \det (A - \lambda I_2) = \det \begin{bmatrix}
					                                                   3 - \lambda & -5          \\
					                                                   -5          & 3 - \lambda
				                                                   \end{bmatrix} \\ &= (3 - \lambda)^2 - 25\\ &= \lambda^2 - 6\lambda - 16.
			      \end{align*}
			      Suas raízes são $\lambda_1 = -2$ e $\lambda_2 = 8$. Tomando $a' = -2$ e $c' = 8$ precisamos resolver o sistema
			      \begin{align*}
				      (A + 2I_2)X = \begin{bmatrix}
					                    0 \\0
				                    \end{bmatrix} \\
				      \begin{bmatrix}
					      5  & -5 \\
					      -5 & 5
				      \end{bmatrix} \begin{bmatrix}
					                    x \\y
				                    \end{bmatrix} = \begin{bmatrix}
					                                    0 \\0
				                                    \end{bmatrix}
			      \end{align*}
			      para encontrar a matriz de rotação. A solução deste sistema é $x = y$ e para determinar a matriz precisamos escolher soluções de norma 1. Assim podemos tomar, por exemplo,
			      \begin{align*}
				      U_1 = (\cos\theta, \sin\theta) = (1/\sqrt{2}, 1/\sqrt{2}) \\
				      U_2 = (-\sin\theta, \cos\theta) = (-1/\sqrt{2}, 1/\sqrt{2}).
			      \end{align*}
			      Daí o ângulo de rotação é $\theta = \arccos(1/\sqrt{2}) = \pi/4$, a matriz de rotação é
			      \[
				      R_\theta = \begin{bmatrix}
					      1/\sqrt{2} & -1/\sqrt{2} \\
					      1/\sqrt{2} & 1/\sqrt{2}
				      \end{bmatrix}
			      \]
			      e
			      \[
				      B = \begin{bmatrix}
					      -2 & 0 \\
					      0  & 8
				      \end{bmatrix}.
			      \]
			      Com isso a equação original torna-se
			      \begin{align*}
				      -2x_1^2 + 8y_1^2 + 8x_1 -16y_1 + 32 = 0                   \\
				      x_1^2 - 4y_1^2 - 4x_1 + 8y_1 -16 = 0                      \\
				      (x_1^2 - 4x_1 + 4) - 4 - 4(y_1^2 - 2y_1 + 1) + 4 - 16 = 0 \\
				      (x_1 - 2)^2 - 4(y_1 - 1)^2 = 16
			      \end{align*}
			      Fazendo a mudança
			      \begin{align*}
				      x_2 = x_1 - 2 \\
				      y_2 = y_1 - 1
			      \end{align*}
			      obtemos a equação
			      \[
				      \dfrac{x_2^2}{16} - \dfrac{y_2^2}{4} = 1
			      \]
			      que é uma hipérbole com focos no eixo $x_2$.
			      \begin{figure}[!h]
				      \centering
				      \caption{Hipérbole $3x^2 +3y^2 - 10xy + 12\sqrt{2}x - 4\sqrt{2}y + 32 = 0$}
				      \includegraphics{hiperbole-rotacionada-transladada.pdf}
			      \end{figure}
		      \end{solucao}
	\end{enumerate}
\end{exemplos}

\begin{teorema}
	Seja $\mathcal{C}$ o conjunto dos pontos do plano que satisfazem a equação
	\[
		ax^2 + bxy + cy^2 + dx + ey + f = 0,
	\]
	com $a$, $b$, $c$, $d$, $e$, $f \in \real$, sendo $a$, $b$ e $c$ não simultaneamente nulos. Sejam $a'$ e $c'$ raízes de
	\[
		p(\lambda) = \det \begin{bmatrix}
			a - \lambda & b/2         \\
			b/2         & c - \lambda
		\end{bmatrix}.
	\]
	\begin{enumerate}[label=({\roman*})]
		\item Temos $a'c' = ac - b^2/4$.
		\item Se $a'c' > 0$, então $\mathcal{C}$ é uma elipse, um ponto ou o conjunto vazio.
		\item Se $a'c' < 0$, então $\mathcal{C}$ é uma hipérbole ou um par de retas concorrentes.
		\item Se $a'c' = 0$, então $\mathcal{C}$ é uma parábola, um par de retas paralelas, uma reta ou o conjunto vazio.
	\end{enumerate}
\end{teorema}
% section rotacoes_de_eixos_metodo_matricial (end)

% section rotacao_e_translacao_de_eixos (end)

% \section{Seções Cônicas} % (fold)
% \label{sec:secoes_conicas}
% \missingfigure{Figuras de seções cônicas}
% section secoes_conicas (end)

% chapter circunferencias_e_conicas (end)
